\documentclass[11pt]{article}
\usepackage[margin=1in]{geometry}
\usepackage{amsmath,amssymb,amsthm}
\usepackage{booktabs}
\usepackage{hyperref}

\newtheorem{theorem}{Theorem}
\newtheorem{definition}{Definition}
\newtheorem{remark}{Remark}

\title{Disproof of Conjecture 00000000427:\\
Spin-Symmetric Plane Partitions in a Cubic Box at $t=1$}
\author{TLMC submission}
\date{September 13, 2026}

\begin{document}
\maketitle

\begin{abstract}
We disprove Conjecture 00000000427, which asserts that plane partitions in a
cubic box that are symmetric under all diagonals and weighted by spin are
counted, at $t=1$, by the closed form
$2^{\lfloor n^2/4\rfloor}\prod_{i=1}^{n}(2i-1)!!/i!$.
At $t=1$ the spin weight is identically $1$, so the claim is purely
combinatorial.  Complete enumeration of the $n=1$ and $n=2$ cases --- under
two natural readings of ``symmetric under all diagonals'' --- shows the true
counts are $2$ and $5$, whereas the conjectured formula gives $1$ and $3$.
The counterexamples are machine-certified in Lean~4 with zero axioms.
\end{abstract}

\section{Definitions}

\begin{definition}[Plane partitions in a cubic box]
A \emph{plane partition in an $n\times n\times n$ box} is an order ideal of
the poset $[n]\times[n]\times[n]$.  Equivalently, writing the cells of the
discrete cube as $\{0,\dots,n-1\}^3$, it is a monotone non-decreasing
$\{0,1\}$-valued height function
\[
h:\{0,\dots,n-1\}^3\to\{0,1\},\qquad
(i\le i')\wedge(j\le j')\wedge(k\le k')\ \Longrightarrow\ h(i,j,k)\le h(i',j',k').
\]
\end{definition}

MacMahon's product formula counts \emph{all} plane partitions in the box:
\[
M(n)\;=\;\prod_{i,j,k=1}^{n}\frac{i+j+k-1}{i+j+k-2},
\qquad M(1)=2,\quad M(2)=20 .
\]

\begin{definition}[Diagonal symmetry]\label{def:sym}
Let $h$ be a height function on the cube.
\begin{itemize}
\item \textbf{Reading A (mirror).} $h$ is \emph{diagonally symmetric} if it is
invariant under every transposition of coordinates:
\[
h(i,j,k)=h(j,i,k)=h(k,j,i)\qquad\text{for all cells }(i,j,k).
\]
\item \textbf{Reading B (cyclic).} $h$ is \emph{diagonally symmetric} if it is
invariant under the cyclic permutation $(i\to j\to k\to i)$:
\[
h(i,j,k)=h(j,k,i)\qquad\text{for all cells }(i,j,k).
\]
\end{itemize}
\end{definition}

On the two-dimensional index set $\{0,1\}^3$ (the $n=2$ case) both readings
generate the same partition of the cube into four orbits
\[
\{(0,0,0)\},\quad
\{(1,0,0),(0,1,0),(0,0,1)\},\quad
\{(1,1,0),(1,0,1),(0,1,1)\},\quad
\{(1,1,1)\},
\]
so a diagonally symmetric height function is exactly one that is constant on
these four orbits, and both readings yield identical counts.  We enumerate
both anyway, so the disproof below is independent of which reading of
``diagonal'' is intended.

\begin{definition}[Spin weight at $t=1$]
The conjecture weights symmetric plane partitions by a ``spin'' parameter
$t$.  At $t=1$ every spin weight equals $1$, so spin does not affect any
count; the conjecture becomes a purely enumerative statement about
Definition~\ref{def:sym}.
\end{definition}

\section{The conjectured formula}

Conjecture 00000000427 asserts
\[
S(n)\;=\;2^{\lfloor n^2/4\rfloor}\prod_{i=1}^{n}\frac{(2i-1)!!}{i!},
\]
where $S(n)$ is the number of diagonally symmetric plane partitions in the
$n\times n\times n$ box at $t=1$.  Evaluating at small $n$:
\[
F(1)=2^{0}\cdot\frac{1!!}{1!}=1,
\qquad
F(2)=2^{1}\cdot\frac{1!!}{1!}\cdot\frac{3!!}{2!}=2\cdot1\cdot\frac{3}{2}=3 .
\]

\section{Enumeration}

We enumerate \emph{all} $\{0,1\}$-valued functions on the cube
($2^{n^3}$ of them: $2$ for $n=1$, $2^8=256$ for $n=2$), keep the monotone
ones, and among those count the diagonally symmetric ones under each reading
of Definition~\ref{def:sym}.  The enumeration is implemented in the
accompanying script \texttt{reproduce.py} and independently re-verified in
Lean~4 (kernel \texttt{decide}, zero axioms; see the \texttt{lean4/}
directory).

\begin{center}
\begin{tabular}{cccccc}
\toprule
$n$ & all functions & monotone & mono.\ + mirror (A) & mono.\ + cyclic (B) & conjectured $F(n)$\\
\midrule
1 & $2$   & $\mathbf{2}$  & $\mathbf{2}$ & $\mathbf{2}$ & $1$\\
2 & $256$ & $\mathbf{20}$ & $\mathbf{5}$ & $\mathbf{5}$ & $3$\\
\bottomrule
\end{tabular}
\end{center}

The monotone counts reproduce MacMahon's totals $M(1)=2$ and $M(2)=20$
exactly, confirming the encoding.  For $n=2$, a diagonally symmetric
monotone height function is determined by its four orbit values
$(a,b,c,d)=\bigl(h(0,0,0),h(1,0,0),h(1,1,0),h(1,1,1)\bigr)$, and monotonicity
forces $a\le b\le c\le d$; the five possibilities are
\[
(0,0,0,0),\ (0,0,0,1),\ (0,0,1,1),\ (0,1,1,1),\ (1,1,1,1),
\]
i.e.\ exactly the non-decreasing $0$--$1$ quadruples.  The $n=1$ cube has a
single cell and exactly two plane partitions (empty and full), both trivially
symmetric.

\section{Verdict}

\begin{theorem}[Disproof of Conjecture 00000000427]
The conjecture is \textbf{FALSE}.  Under either reading of ``symmetric under
all diagonals'':
\begin{align*}
n=1:&\quad S(1)=2\neq 1=F(1),\\
n=2:&\quad S(2)=5\neq 3=F(2).
\end{align*}
\end{theorem}

\begin{remark}[Scope]
MacMahon's product formula itself is correct and is not contested: our
enumeration reproduces $M(1)=2$, $M(2)=20$.  What is refuted is the claimed
closed form for the \emph{diagonal-symmetry} count at $t=1$, together with the
claimed shape of its ratio to the unsymmetrized MacMahon formula (e.g.\ at
$n=2$ the true ratio is $5/20=1/4$, not $3/20$).
\end{remark}

\begin{remark}[Robustness of the encoding]
Under the alternative, non-standard convention in which the height function is
$3$-valued, $h:\{0,1\}^3\to\{0,1,2\}$, the monotone functions number $168$
(= plane partitions in a $2\times2\times2\times2$ box by MacMahon's formula)
and the diagonally symmetric ones number $15$ under both readings --- again
different from the conjectured value $3$.  The disproof is thus independent of
that encoding choice as well.
\end{remark}

\section{Machine certification}

The \texttt{lean4/} directory contains a Lean~4 project
(toolchain \texttt{leanprover/lean4:v4.33.1}, library \texttt{Main}) proving:
\begin{itemize}
\item \texttt{macMahon2}: filtering the $2^8=256$ height functions of the
  $2\times2\times2$ cube by monotonicity leaves exactly $20$;
\item \texttt{symCount2}: monotone $+$ mirror-symmetric count $= 5$;
\item \texttt{cycSymCount2}: monotone $+$ cyclic-symmetric count $= 5$;
\item \texttt{formula2}: the conjectured value $2^1\cdot(1!!/1!)\cdot(3!!/2!)=3$
  (stated as the exact cross-multiplied computation
  $2^{\lfloor 4/4\rfloor}\cdot(1\cdot3)/(1\cdot2)=3$ over $\mathbb{N}$, since
  core Lean's $\mathbb{Q}$ operations do not kernel-reduce);
\item \texttt{refute2}: the conjunction of the above with
  $\texttt{symCount2}\neq 3$ and $\texttt{cycSymCount2}\neq 3$;
\item \texttt{macMahon1}, \texttt{symCount1}, \texttt{formula1},
  \texttt{refute1}: the $n=1$ counterexample $2\neq 1$.
\end{itemize}
All proofs are by kernel \texttt{decide}/\texttt{rfl}; \texttt{Check.lean}
prints \texttt{\#print axioms} for every theorem, each of which reports that it
depends on \emph{no axioms}.  There is no \texttt{sorry} anywhere.

\end{document}
